El dominio de la terminal es fundamental para un desarrollador porque optimiza el tiempo y aumenta la productividad al gestionar proyectos. A diferencia de la interfaz gráfica, la línea de comandos permite navegar por estructuras de archivos complejas de forma rápida, automatizar tareas repetitivas y ejecutar acciones masivas con pocos comandos. Además, la mayoría de las herramientas esenciales en el desarrollo moderno, como los sistemas de control de versiones como Git, los gestores de dependencias y los servidores en la nube, operan de forma nativa o más eficiente a través de la terminal. Conocer estos comandos básicos no solo brinda un control absoluto sobre el entorno de trabajo local, sino que también es un requisito indispensable para administrar servidores remotos donde no existe una interfaz visual. En conclusión, la terminal es el puente principal entre el código y el sistema operativo, haciendo al desarrollador más ágil y versátil.